%% notes-tex/database/ecoli-datasets/ecoli-datasets.tex
%% [[experiments.database.ecoli-datasets]]
%% https://github.com/Mjvolk3/torchcell/tree/main/notes-tex/database/ecoli-datasets/ecoli-datasets.tex

\documentclass[10pt,a4paper]{article}
\usepackage[draft]{tcdoc}
\usepackage{float}
\usepackage{pdflscape}

%% Every number the prose quotes is a macro generated from the snapshot.
%% GENERATED FILE -- do not hand-edit.
%% SOURCE: experiments/database/scripts/build_ecoli_genslm_benchmark_table.py
%% Every number the prose quotes, as a macro, so it is read from the snapshot.
\newcommand{\snapshotDate}{2026-10-21}
\newcommand{\ecoliDatasetCount}{68}
\newcommand{\ecoliInstanceTotal}{4,561,865}
\newcommand{\ecoliSectionCount}{13}
\newcommand{\ecoliBestCount}{35}
\newcommand{\ecoliGoodCount}{29}
\newcommand{\ecoliWeakCount}{1}
\newcommand{\ecoliNoneCount}{3}
\newcommand{\pickCount}{8}
\newcommand{\coreCount}{4}
\newcommand{\coreGenesMin}{2,359}
\newcommand{\coreGenesMax}{4,181}
\newcommand{\coreLabelsTotal}{1,386,311}
\newcommand{\priceEnv}{104}
\newcommand{\goodallGenes}{4,181}
\newcommand{\goodallLabels}{8,309}
\newcommand{\goodallEnv}{2}
\newcommand{\goodallInstances}{8,309}
\newcommand{\goodallIdentifier}{BW25113 locus tag, released as gene symbol}
\newcommand{\nicholsGenes}{3,949}
\newcommand{\nicholsLabels}{1,286,962}
\newcommand{\nicholsEnv}{324}
\newcommand{\nicholsInstances}{1,286,962}
\newcommand{\nicholsIdentifier}{BW25113 locus tag, released as ECK accession}
\newcommand{\schmidtGenes}{2,359}
\newcommand{\schmidtLabels}{48,335}
\newcommand{\schmidtEnv}{21}
\newcommand{\schmidtInstances}{21}
\newcommand{\schmidtIdentifier}{gene symbol (readout key)}
\newcommand{\babuGenes}{3,880}
\newcommand{\babuLabels}{42,705}
\newcommand{\babuEnv}{1}
\newcommand{\babuInstances}{42,705}
\newcommand{\babuIdentifier}{MG1655 b-number, released as Keio JW number; b-number with a Keio sub-number; gene symbol}
\newcommand{\liGenes}{3,410}
\newcommand{\liLabels}{8,644}
\newcommand{\liEnv}{3}
\newcommand{\liInstances}{3}
\newcommand{\liIdentifier}{gene symbol (readout key)}
\newcommand{\preciseGenes}{4,257}
\newcommand{\preciseLabels}{4,405,995}
\newcommand{\preciseEnv}{119}
\newcommand{\preciseInstances}{1,035}
\newcommand{\preciseIdentifier}{b-number (readout key)}
\newcommand{\camposGenes}{3,664}
\newcommand{\camposLabels}{53,508}
\newcommand{\camposEnv}{1}
\newcommand{\camposInstances}{3,664}
\newcommand{\camposIdentifier}{BW25113 locus tag}
\newcommand{\wangGenes}{4,223}
\newcommand{\wangLabels}{20,158}
\newcommand{\wangEnv}{5}
\newcommand{\wangInstances}{20,158}
\newcommand{\wangIdentifier}{MG1655 b-number}
\newcommand{\guideStoreCount}{6}


\title{\vspace{-1.5em}The \emph{E. coli} Datasets, and the Short List That Can
Benchmark a Gene-Sequence Encoding}
\author{Michael Volk}
\date{\today}

\begin{document}
\maketitle
\thispagestyle{empty}

\noindent
The dev tree serves \ecoliDatasetCount{} \org{Escherichia coli} datasets holding
\ecoliInstanceTotal{} records, built as LMDBs from their own releases and read here from
the snapshot of \snapshotDate{} (Table~\ref{tab:ecoli-datasets}). Every column of that
table is read from a built store; nothing in it is typed by hand.

A \term{gene-sequence encoding} is one vector per coding sequence and nothing else: no
cell graph, no annotation, no gene identity beyond what the sequence carries. GenSLM is
the instance at hand, a codon language model pretrained on prokaryotic genes, and
\file{torchcell/datasets/genslm.py} already builds such a store for any genome whose
genes expose a coding sequence. Benchmarking that encoding means asking, dataset by
dataset, whether the vector of a gene the model has not been fit on predicts that gene's
label. Of the \ecoliDatasetCount{} datasets, \ecoliBestCount{} attach a label to one gene,
one gene pair or one inserted construct (Table~\ref{tab:tiers}); far fewer carry enough
labeled genes, on a task worth scoring, to make a benchmark. Sec.~\ref{sec:benchmark}
sets the rule and gives the short list (Table~\ref{tab:shortlist}): \coreCount{} core
datasets, one per task family, each with between \coreGenesMin{} and \coreGenesMax{}
labeled genes, and a second tier of the same size. Sec.~\ref{sec:caveats} names what a
score on any of them does not show.

\vspace{0.5em}
\tccontents
\clearpage

%% notes-tex/database/ecoli-datasets/sections/datasets.tex
%% [[experiments.database.ecoli-datasets]]
%% https://github.com/Mjvolk3/torchcell/tree/main/notes-tex/database/ecoli-datasets/sections/datasets.tex
%% SOURCE: hand-authored prose; every number quoted here is a macro from
%% tables/numbers.tex, generated by
%% experiments/database/scripts/build_ecoli_genslm_benchmark_table.py off the snapshot
%% written by experiments/database/scripts/build_bacteria_datasets_table.py

\section{The \emph{E. coli} datasets\secstatus{todo}}\label{sec:datasets}

Table~\ref{tab:ecoli-datasets} lists the \ecoliDatasetCount{} datasets in
\ecoliSectionCount{} phenotype families. \term{Genotypes} and \term{Env} are the
distinct stored genotype and environment payloads, \term{Instances} the record count,
and \term{Signal} the gzip size in bits of the concatenated stored instances, the same
relative proxy the yeast table uses. A dataset whose strains are all the one wild type,
such as a proteome panel across growth conditions, shows one genotype and as many
environments as conditions; a deletion screen shows one genotype per strain.

Two columns are specific to this document. The \term{sequence-only fit} tier says how
well a model that sees only sequence embeddings fits the dataset, and the \term{label
unit} names what one label attaches to once the cell graph is dropped. Both are derived
by the probe script from where the label sits in the cell graph (a whole-strain readout,
a node, a gene pair, a metabolite) and from the modal number of edits in a sample of
400 genotypes per store. The tiers are structural: \emph{best} when the label sits on
one gene, one gene pair or triple, or one inserted construct, so that a per-gene or
per-pair representation is the whole input; \emph{good} when the label needs the readout
entity crossed with the perturbed genes or with the condition, which a pair of vectors
still carries; \emph{weak} for a multi-edit strain where only its heterologous genes have
a sequence to stand on; \emph{no} for a strain readout with no single edit to attach a
sequence to. Table~\ref{tab:tiers} sums the datasets and instances per tier. A tier is a
reading of the genotype, not a model result; it says what a sequence-only model could
be given, not what it would score.

%% GENERATED FILE -- do not hand-edit.
%% SOURCE: experiments/database/scripts/build_ecoli_genslm_benchmark_table.py
%% off experiments/database/results/pre-build/2026-10-21/ecoli_datasets.json
\begin{table}[H]
\centering
\footnotesize
\setlength{\tabcolsep}{4pt}
\caption{\emph{E. coli} datasets and instances per sequence-only fit tier, as \texttt{sequence\_fit} in the probe script derives it from where the label sits and the modal edit count of the genotype sample.}
\label{tab:tiers}
\begin{tabular}{@{}l r r >{\raggedright\arraybackslash}p{88mm}@{}}
\toprule
Sequence-only fit & Datasets & Instances & What the tier means \\
\midrule
best & 35 & 639,394 & the label sits on one gene, a gene pair or triple, or one inserted construct \\
good & 29 & 3,921,964 & readout entity x perturbed gene(s), or gene x condition \\
weak & 1 & 72 & a multi-edit strain, only its heterologous genes carry a sequence \\
no & 3 & 435 & a natural genome, or a whole-strain readout with no single edit \\
\midrule
\textbf{Total} & \textbf{68} & \textbf{4,561,865} &  \\
\bottomrule
\end{tabular}
\end{table}


\clearpage
%% GENERATED FILE -- do not hand-edit.
%% SOURCE: experiments/database/scripts/build_ecoli_genslm_benchmark_table.py
%% off experiments/database/results/pre-build/2026-10-21/ecoli_datasets.json
\begin{landscape}
\begingroup
\scriptsize
\setlength{\tabcolsep}{3pt}
\begin{longtable}{@{}l r r r >{\raggedright\arraybackslash}p{34mm} l l r l >{\raggedright\arraybackslash}p{40mm}@{}}
\caption{The \emph{E. coli} datasets served from the dev tree, snapshot 2026-10-21, grouped by phenotype family. \emph{Genotypes} and \emph{Env} count distinct stored genotype and environment payloads; \emph{Instances} is the record count; \emph{Signal} the gzip size in bits of the concatenated stored instances (perturbation, environment and phenotype). \emph{Seq.\ fit} is the sequence-only fit tier of Table~\ref{tab:tiers} and \emph{Label unit} what one label attaches to without the cell graph, both read from the genotype of 400 evenly spaced records per store.}\label{tab:ecoli-datasets}\\
\toprule
Dataset & Genotypes & Env & Instances & Phenotype & Shape & Role & Signal (bits) & Seq. fit & Label unit \\
\midrule
\endfirsthead
\multicolumn{10}{@{}l}{\footnotesize\emph{Table~\ref{tab:ecoli-datasets}, continued}}\\
\toprule
Dataset & Genotypes & Env & Instances & Phenotype & Shape & Role & Signal (bits) & Seq. fit & Label unit \\
\midrule
\endhead
\multicolumn{10}{@{}l}{\textbf{Environmental / chemogenomic response}} \\
CrisprPineneToleranceNiu2019 & 54 & 1 & 54 & environment\_\allowbreak{}response (log2\_\allowbreak{}ratio) & scalar & global & $4.1\times10^{4}$ & best & 1 gene (1-6 edits) \\
CrispriChemgenChoe2025 & 39,291 & 12 & 470,502 & environment\_\allowbreak{}response (log2\_\allowbreak{}ratio) & scalar & global & $2.4\times10^{8}$ & good & 1 gene x condition \\
CrispriGeneFitnessWang2018 & 13,641 & 5 & 20,158 & environment\_\allowbreak{}response (log2\_\allowbreak{}ratio) & scalar & global & $1.5\times10^{7}$ & good & 1 gene x condition (1-7 edits) \\
CrispriGuideFfaEnrichmentFang2025 & 15,711 & 1 & 15,711 & environment\_\allowbreak{}response (log2\_\allowbreak{}ratio) & scalar & global & $7.3\times10^{6}$ & best & 1 gene (1-5 edits) \\
CrispriGuideFitnessWang2018 & 53,203 & 5 & 240,481 & environment\_\allowbreak{}response (log2\_\allowbreak{}ratio) & scalar & global & $1.6\times10^{8}$ & good & 1 gene x condition \\
CrispriKnockdownCui2018 & 71,696 & 1 & 143,392 & environment\_\allowbreak{}response (log2\_\allowbreak{}ratio) & scalar & global & $7.3\times10^{7}$ & best & 1 gene (1-7 edits) \\
CrispriScreenRousset2018 & 40,592 & 4 & 92,252 & environment\_\allowbreak{}response (log2\_\allowbreak{}ratio) & scalar & global & $4.2\times10^{7}$ & good & 1 gene x condition \\
DoublingTimeCaglar2017 & 1 & 16 & 55 & environment\_\allowbreak{}response (growth\_\allowbreak{}rate) & scalar & global & $3.2\times10^{4}$ & no & unperturbed strain readout x condition \\
EnvChemgenGirgis2009 & 3,976 & 17 & 66,287 & environment\_\allowbreak{}response (z\_\allowbreak{}score) & scalar & global & $3.7\times10^{7}$ & good & 1 gene x condition \\
EnvChemgenNichols2011 & 3,967 & 324 & 1,286,962 & environment\_\allowbreak{}response (z\_\allowbreak{}score) & scalar & global & $4.1\times10^{8}$ & good & 1 gene x condition \\
EnvChemgenShiver2016 & 3,846 & 51 & 210,998 & environment\_\allowbreak{}response (z\_\allowbreak{}score) & scalar & global & $9.9\times10^{7}$ & good & 1 gene x condition \\
EnvChemgenWang2015 & 46 & 1 & 46 & environment\_\allowbreak{}response (log2\_\allowbreak{}ratio) & scalar & global & $3.0\times10^{4}$ & best & 1 gene \\
EnvChemgenWang2024 & 4,419 & 6 & 26,514 & environment\_\allowbreak{}response (log2\_\allowbreak{}ratio) & scalar & global & $9.3\times10^{6}$ & good & 1 gene x condition \\
FitnessBrowserEcoli & 3,789 & 3 & 22,734 & environment\_\allowbreak{}response (log2\_\allowbreak{}ratio) & scalar & global & $1.1\times10^{7}$ & good & 1 gene x condition \\
GrowthRateLamoureux2023 & 21 & 54 & 354 & environment\_\allowbreak{}response (growth\_\allowbreak{}rate) & scalar & global & $2.1\times10^{5}$ & no & unperturbed strain readout x condition (0-1 edits) \\
GrowthRateS23Schmidt2016 & 1 & 22 & 26 & environment\_\allowbreak{}response (growth\_\allowbreak{}rate) & scalar & global & $2.0\times10^{4}$ & no & unperturbed strain readout x condition \\
GrowthYieldBaba2006 & 3,912 & 3 & 11,735 & environment\_\allowbreak{}response (optical\_\allowbreak{}density) & scalar & global & $3.3\times10^{6}$ & good & 1 gene x condition \\
MismatchCrispriFitnessHawkins2020 & 24,186 & 1 & 24,186 & environment\_\allowbreak{}response (relative\_\allowbreak{}growth\_\allowbreak{}rate) & scalar & global & $1.4\times10^{7}$ & best & 1 gene \\
PhageRbTnseqMutalik2020 & 3,716 & 71 & 287,815 & environment\_\allowbreak{}response (log2\_\allowbreak{}ratio) & scalar & global & $2.1\times10^{8}$ & good & 1 gene x condition \\
RbTnseqPrice2018Ecoli & 3,789 & 104 & 598,662 & environment\_\allowbreak{}response (log2\_\allowbreak{}ratio) & scalar & global & $4.0\times10^{8}$ & good & 1 gene x condition \\
RbTnseqTeteneva2024 & 3,691 & 18 & 66,162 & environment\_\allowbreak{}response (log2\_\allowbreak{}ratio) & scalar & global & $2.7\times10^{7}$ & good & 1 gene x condition \\
RbTnseqWetmore2015 & 3,646 & 85 & 353,662 & environment\_\allowbreak{}response (log2\_\allowbreak{}ratio) & scalar & global & $2.4\times10^{8}$ & good & 1 gene x condition \\
\addlinespace
\multicolumn{10}{@{}l}{\textbf{Expression (RNA-seq)}} \\
MrnaFractionBalakrishnan2022 & 3 & 14 & 28 & mrna\_\allowbreak{}number\_\allowbreak{}fraction (rnaseq\_\allowbreak{}mrna\_\allowbreak{}number\_\allowbreak{}fraction) & vector (4342) & node & $7.2\times10^{6}$ & best & 1 gene, unperturbed strain x condition (0-2 edits) \\
RnaseqCaglar2017 & 1 & 57 & 152 & expression\_\allowbreak{}tpm (rnaseq\_\allowbreak{}tpm) & vector (4196) & node & $6.7\times10^{7}$ & best & 1 gene, unperturbed strain x condition \\
RnaseqGuzman2018 & 7 & 2 & 14 & expression\_\allowbreak{}tpm (rnaseq\_\allowbreak{}tpm) & vector (4329) & node & $4.5\times10^{6}$ & good & readout gene x 3 perturbed gene(s) x condition (0-5 edits) \\
RnaseqLamoureux2023 & 74 & 119 & 1,035 & expression\_\allowbreak{}tpm (rnaseq\_\allowbreak{}tpm) & vector (4257) & node & $5.8\times10^{8}$ & best & 1 gene, unperturbed strain x condition (0-3 edits) \\
RnaseqPublicK12Lamoureux2023 & 72 & 198 & 1,675 & expression\_\allowbreak{}tpm (rnaseq\_\allowbreak{}tpm\_\allowbreak{}from\_\allowbreak{}released\_\allowbreak{}counts) & vector (4355) & node & $9.5\times10^{8}$ & best & 1 gene, unperturbed strain x condition (0-4 edits) \\
\addlinespace
\multicolumn{10}{@{}l}{\textbf{Fitness}} \\
CarbonSourceTong2020 & 3,714 & 30 & 111,420 & fitness & scalar & global & $4.8\times10^{7}$ & good & 1 gene x condition \\
CrispriCrossRachwalski2024 & 14,209 & 15 & 48,900 & fitness & scalar & global & $1.9\times10^{7}$ & good & gene pair, strain readout x condition (1-2 edits) \\
GrowthAucRapp2026 & 1,515 & 1 & 1,515 & fitness & scalar & global & $7.6\times10^{5}$ & best & 1 gene \\
GrowthRateCampos2018 & 3,664 & 1 & 3,664 & fitness & scalar & global & $1.1\times10^{6}$ & best & 1 gene \\
GrowthRateChoe2019 & 6 & 1 & 8 & fitness & scalar & global & $1.1\times10^{4}$ & best & 1 gene (0-21 edits) \\
GrowthRateSchmidt2016 & 3 & 2 & 6 & fitness & scalar & global & $8.5\times10^{3}$ & good & 1 gene x condition \\
GrowthWang2015 & 46 & 1 & 46 & fitness & scalar & global & $2.2\times10^{4}$ & best & 1 gene \\
TranscriptionFactorKnockoutChoe2019 & 2 & 1 & 2 & fitness & scalar & global & $5.8\times10^{3}$ & best & 1 gene \\
\addlinespace
\multicolumn{10}{@{}l}{\textbf{Flux}} \\
FluxIshii2007 & 25 & 5 & 28 & net\_\allowbreak{}flux (c13\_\allowbreak{}mfa\_\allowbreak{}net\_\allowbreak{}flux\_\allowbreak{}percent\_\allowbreak{}of\_\allowbreak{}specific\_\allowbreak{}glucose\_\allowbreak{}uptake) & vector (42) & bipartite node & $6.0\times10^{4}$ & good & reaction x 1 perturbed gene(s) x condition (0-1 edits) \\
FluxLi2021 & 3 & 1 & 3 & net\_\allowbreak{}flux (c13\_\allowbreak{}mfa\_\allowbreak{}net\_\allowbreak{}flux\_\allowbreak{}mmol\_\allowbreak{}per\_\allowbreak{}gdw\_\allowbreak{}per\_\allowbreak{}h) & vector (99) & bipartite node & $5.7\times10^{4}$ & good & reaction x 1 perturbed gene(s) \\
\addlinespace
\multicolumn{10}{@{}l}{\textbf{Gene interaction}} \\
GeneInteractionBabu2014 & 42,703 & 1 & 42,705 & gene\_\allowbreak{}interaction & scalar & edge & $1.2\times10^{7}$ & best & gene pair \\
GeneInteractionButland2008 & 314,769 & 1 & 314,769 & gene\_\allowbreak{}interaction & scalar & edge & $8.2\times10^{7}$ & best & gene pair \\
\addlinespace
\multicolumn{10}{@{}l}{\textbf{Metabolite}} \\
ExometaboliteBrunk2016 & 9 & 28 & 126 & metabolite\_\allowbreak{}level (hplc\_\allowbreak{}extracellular\_\allowbreak{}concentration\_\allowbreak{}g\_\allowbreak{}per\_\allowbreak{}l) & vector (6) & bipartite node & $2.3\times10^{5}$ & good & metabolite x 9 perturbed gene(s) x condition (0-10 edits) \\
MetaboliteIntensityRapp2026 & 407 & 1 & 407 & metabolite\_\allowbreak{}level (fi\_\allowbreak{}ms\_\allowbreak{}annotated\_\allowbreak{}feature\_\allowbreak{}absolute\_\allowbreak{}intensity\_\allowbreak{}mean\_\allowbreak{}of\_\allowbreak{}2\_\allowbreak{}plates) & vector (2) & bipartite node & $3.9\times10^{5}$ & good & metabolite x 1 perturbed gene(s) \\
MetabolomeBrunk2016 & 9 & 26 & 117 & metabolite\_\allowbreak{}level (lc\_\allowbreak{}ms\_\allowbreak{}concentration\_\allowbreak{}um) & vector (51) & bipartite node & $4.5\times10^{5}$ & good & metabolite x 9 perturbed gene(s) x condition (0-10 edits) \\
MetabolomeFuhrer2017 & 3,806 & 1 & 3,806 & metabolite\_\allowbreak{}level (fia\_\allowbreak{}tof\_\allowbreak{}ms\_\allowbreak{}ion\_\allowbreak{}modified\_\allowbreak{}z\_\allowbreak{}score) & vector (7534) & bipartite node & $1.9\times10^{9}$ & good & metabolite x 1 perturbed gene(s) \\
MetabolomeIshii2007 & 25 & 5 & 28 & metabolite\_\allowbreak{}level (ce\_\allowbreak{}tofms\_\allowbreak{}intracellular\_\allowbreak{}concentration\_\allowbreak{}mm) & vector (130) & bipartite node & $4.9\times10^{5}$ & good & metabolite x 1 perturbed gene(s) x condition (0-1 edits) \\
MetabolomeRapp2026 & 1,497 & 1 & 1,497 & metabolite\_\allowbreak{}level (fi\_\allowbreak{}ms\_\allowbreak{}iml1515\_\allowbreak{}feature\_\allowbreak{}fold\_\allowbreak{}change\_\allowbreak{}vs\_\allowbreak{}batch\_\allowbreak{}median\_\allowbreak{}mean\_\allowbreak{}of\_\allowbreak{}2\_\allowbreak{}plates) & vector (1321) & bipartite node & $5.6\times10^{8}$ & good & metabolite x 1 perturbed gene(s) \\
MetabolomeSchastnaya2021 & 103 & 5 & 200 & metabolite\_\allowbreak{}level (fia\_\allowbreak{}tof\_\allowbreak{}ms\_\allowbreak{}ion\_\allowbreak{}log2\_\allowbreak{}fold\_\allowbreak{}change\_\allowbreak{}vs\_\allowbreak{}wild\_\allowbreak{}type) & vector (284) & bipartite node & $3.6\times10^{6}$ & good & metabolite x 1 perturbed gene(s) x condition \\
TargetedMetabolomeRapp2026 & 407 & 1 & 407 & metabolite\_\allowbreak{}level (lc\_\allowbreak{}msms\_\allowbreak{}eic\_\allowbreak{}peak\_\allowbreak{}height\_\allowbreak{}fold\_\allowbreak{}change\_\allowbreak{}vs\_\allowbreak{}library\_\allowbreak{}median) & vector (2) & bipartite node & $3.6\times10^{5}$ & good & metabolite x 1 perturbed gene(s) \\
\addlinespace
\multicolumn{10}{@{}l}{\textbf{Morphology}} \\
MorphologyCampos2018 & 3,664 & 1 & 3,664 & morphology & vector (15) & global & $8.5\times10^{6}$ & best & 1 gene \\
\addlinespace
\multicolumn{10}{@{}l}{\textbf{Phosphoproteome}} \\
PhosphositeIntensityPotel2018 & 1 & 2 & 12 & phosphosite\_\allowbreak{}intensity (fe3\_\allowbreak{}imac\_\allowbreak{}lc\_\allowbreak{}msms\_\allowbreak{}maxquant\_\allowbreak{}summed\_\allowbreak{}xic) & vector (989) & node & $1.3\times10^{7}$ & best & 1 gene, unperturbed strain x condition \\
\addlinespace
\multicolumn{10}{@{}l}{\textbf{Product titer}} \\
BiofuelTiterBrunk2016 & 8 & 18 & 72 & titer & scalar & global & $1.5\times10^{5}$ & weak & set of 9 edits incl. heterologous genes (7-10 edits) \\
IsopentenolTiterFoo2014 & 9 & 1 & 9 & titer & scalar & global & $1.7\times10^{4}$ & best & 1 heterologous construct \\
PineneFoldChangeNiu2019 & 40 & 1 & 40 & titer & scalar & global & $2.3\times10^{4}$ & best & 1 gene (1-6 edits) \\
\addlinespace
\multicolumn{10}{@{}l}{\textbf{Promoter activity}} \\
PromoterReporterMohiuddin2022 & 1,809 & 27 & 69,480 & promoter\_\allowbreak{}activity & scalar & node & $3.3\times10^{7}$ & best & inserted DNA (promoter or reporter) x condition \\
\addlinespace
\multicolumn{10}{@{}l}{\textbf{Protein abundance}} \\
ProteinFoldChangeCaglar2017 & 1 & 8 & 8 & protein\_\allowbreak{}fold\_\allowbreak{}change (deseq2\_\allowbreak{}wald\_\allowbreak{}log2\_\allowbreak{}fold\_\allowbreak{}change\_\allowbreak{}design\_\allowbreak{}batch\_\allowbreak{}plus\_\allowbreak{}condition) & vector (4180) & node & $1.1\times10^{7}$ & best & 1 gene, unperturbed strain x condition \\
ProteomeBrunk2016 & 9 & 18 & 81 & protein\_\allowbreak{}abundance (srm\_\allowbreak{}protein\_\allowbreak{}peak\_\allowbreak{}area) & vector (68) & node & $4.2\times10^{5}$ & good & readout gene x 9 perturbed gene(s) x condition (0-10 edits) \\
ProteomeCaglar2017 & 1 & 40 & 105 & protein\_\allowbreak{}abundance (lcmsms\_\allowbreak{}spectral\_\allowbreak{}count\_\allowbreak{}deseq2\_\allowbreak{}size\_\allowbreak{}factor\_\allowbreak{}normalized) & vector (4196) & node & $2.5\times10^{7}$ & best & 1 gene, unperturbed strain x condition \\
ProteomeIshii2007 & 25 & 5 & 28 & protein\_\allowbreak{}abundance (lc\_\allowbreak{}ms\_\allowbreak{}ms\_\allowbreak{}absolute\_\allowbreak{}mg\_\allowbreak{}protein\_\allowbreak{}per\_\allowbreak{}g\_\allowbreak{}dry\_\allowbreak{}cell\_\allowbreak{}weight) & vector (57) & node & $2.9\times10^{5}$ & best & 1 gene, its own label x condition (0-1 edits) \\
ProteomeMori2021 & 1 & 1 & 7 & protein\_\allowbreak{}abundance (absolute\_\allowbreak{}protein\_\allowbreak{}mass\_\allowbreak{}fraction\_\allowbreak{}xtop\_\allowbreak{}dia\_\allowbreak{}swath\_\allowbreak{}riboprofiling\_\allowbreak{}scaled) & vector (1897) & node & $1.1\times10^{6}$ & best & 1 gene, unperturbed strain \\
ProteomeSchmidt2016 & 1 & 21 & 21 & protein\_\allowbreak{}abundance (absolute\_\allowbreak{}protein\_\allowbreak{}copies\_\allowbreak{}per\_\allowbreak{}cell\_\allowbreak{}label\_\allowbreak{}free\_\allowbreak{}sid\_\allowbreak{}anchored) & vector (2359) & node & $1.1\times10^{7}$ & best & 1 gene, unperturbed strain x condition \\
ProteomeSrmSet1Schmidt2016 & 1 & 18 & 18 & protein\_\allowbreak{}abundance (absolute\_\allowbreak{}protein\_\allowbreak{}copies\_\allowbreak{}per\_\allowbreak{}cell\_\allowbreak{}srm\_\allowbreak{}sid\_\allowbreak{}dataset1\_\allowbreak{}technical\_\allowbreak{}sd) & vector (41) & node & $1.5\times10^{5}$ & best & 1 gene, unperturbed strain x condition \\
ProteomeSrmSet2Schmidt2016 & 1 & 21 & 21 & protein\_\allowbreak{}abundance (absolute\_\allowbreak{}protein\_\allowbreak{}copies\_\allowbreak{}per\_\allowbreak{}cell\_\allowbreak{}srm\_\allowbreak{}sid\_\allowbreak{}dataset2\_\allowbreak{}biological\_\allowbreak{}sd) & vector (31) & node & $1.3\times10^{5}$ & best & 1 gene, unperturbed strain x condition \\
\addlinespace
\multicolumn{10}{@{}l}{\textbf{Protein synthesis / turnover}} \\
ProteinSynthesisRateLi2014 & 1 & 3 & 3 & synthesis\_\allowbreak{}rate (ribosome\_\allowbreak{}profiling\_\allowbreak{}footprint\_\allowbreak{}density\_\allowbreak{}absolute\_\allowbreak{}synthesis\_\allowbreak{}rate) & vector (3041) & node & $3.7\times10^{5}$ & best & 1 gene, unperturbed strain x condition \\
ProteinTurnoverGupta2024 & 6 & 8 & 13 & degradation\_\allowbreak{}rate (n15\_\allowbreak{}ammonium\_\allowbreak{}tmtproc\_\allowbreak{}total\_\allowbreak{}turnover\_\allowbreak{}rate\_\allowbreak{}per\_\allowbreak{}hour) & vector (2553) & node & $1.1\times10^{7}$ & best & 1 gene, unperturbed strain x condition (0-3 edits) \\
\addlinespace
\multicolumn{10}{@{}l}{\textbf{Viability}} \\
GeneEssentialityBaba2006 & 4,288 & 1 & 4,288 & is\_\allowbreak{}essential & scalar & node & $1.8\times10^{6}$ & best & 1 gene, its own label \\
GeneEssentialityGerdes2003 & 3,739 & 1 & 3,739 & is\_\allowbreak{}essential & scalar & node & $2.0\times10^{6}$ & best & 1 gene, its own label \\
GeneEssentialityGoodall2018 & 8,309 & 2 & 8,309 & is\_\allowbreak{}essential & scalar & node & $5.2\times10^{6}$ & best & 1 gene, its own label x condition \\
GeneEssentialityPrice2018Ecoli & 324 & 1 & 324 & is\_\allowbreak{}essential & scalar & node & $2.6\times10^{5}$ & best & 1 gene, its own label \\
GradedGeneEssentialityGoodall2018 & 317 & 2 & 317 & essentiality\_\allowbreak{}call & scalar & node & $2.2\times10^{5}$ & best & 1 gene, its own label x condition \\
\midrule
\textbf{Total (68 datasets)} &  &  & \textbf{4,561,865} &  &  &  & \textbf{$6.4\times10^{9}$} &  &  \\
\bottomrule
\end{longtable}
\endgroup
\end{landscape}


%% notes-tex/database/ecoli-datasets/sections/benchmark.tex
%% [[experiments.database.ecoli-datasets]]
%% https://github.com/Mjvolk3/torchcell/tree/main/notes-tex/database/ecoli-datasets/sections/benchmark.tex
%% SOURCE: hand-authored prose; every number quoted here is a macro from
%% tables/numbers.tex, generated by
%% experiments/database/scripts/build_ecoli_genslm_benchmark_table.py, and the pick
%% list is the PICKS registry in that script

\section{The short list for a gene-sequence encoding benchmark\secstatus{todo}}\label{sec:benchmark}

\subsection{What is being benchmarked\secstatus{todo}}

GenSLM (Zvyagin et al.\ 2023, mirror key \texttt{zvyaginGenSLMsGenomescaleLanguage2023})
tokenizes a gene as codons and was pretrained on more than 110 million prokaryotic gene
sequences from BV-BRC. \file{torchcell/datasets/genslm.py} embeds each gene's spliced
coding sequence, in frame from the start codon and truncated to the first 2,048 codons
when longer, and stores the masked mean of the last hidden layer: 512 values per gene
for the 25M checkpoint and 1,840 for the 250M one, both of which are on GilaHyper with
a sha256 manifest. The stores built so far are for \org{P. putida} KT2440; the
\org{E. coli} K-12 MG1655 and BW25113 genomes expose a coding sequence through the same
protocol, so their stores are a build away, not a code change.

The question a benchmark asks of such an encoding is narrow. Take a dataset whose label
attaches to a gene, hold out a set of genes, fit a small head on the vectors of the
remaining genes, and score the held-out ones. That score has meaning only against two
controls at the same width: a one-hot or random vector per gene, which carries gene
identity and nothing else, and a codon-frequency vector, which carries the one feature
of a coding sequence a language model is sure to have learned. A codon model earns its
place when it beats both; it is doing something different from either when it beats
one and not the other.

\subsection{The selection rule\secstatus{todo}}

Three conditions, applied in order to every row of Table~\ref{tab:ecoli-datasets}.

\pfstep{The label sits on one gene or one gene pair} The sequence-only fit tier must be
\emph{best}, or \emph{good} only because the gene is crossed with a condition, since the
condition is a second input the gene encoding does not have to carry. This removes
every whole-strain readout: the titers, the growth rates of an engineered strain, the
multi-omics panels of a wild type where the label is a profile of a strain rather than
of a gene. It also removes the \guideStoreCount{} stores whose label sits on a guide
RNA rather than on the gene it targets (Sec.~\ref{sec:caveats}).

\pfstep{Enough labeled genes for a held-out split} A benchmark scored on a few dozen
held-out genes cannot separate two encodings. Every core pick carries between
\coreGenesMin{} and \coreGenesMax{} genes with a label, read from a full scan of the
store, so a ten percent test fold holds hundreds of genes. The count that matters is
genes, not instances: a proteome of \schmidtEnv{} conditions is \schmidtInstances{}
instances and \schmidtGenes{} labeled genes, and it is the latter that is split.

\pfstep{One primary per task family} A short list is one dataset per kind of question,
with the datasets that pose the same question named beside it rather than listed on
their own. The families are a binary call per gene, a quantitative fitness per gene
and condition, a molecular level of the gene's own product, and a pairwise label.

\subsection{The picks\secstatus{todo}}

Table~\ref{tab:shortlist} is the list. The four core rows cover the four families and
should be run first; the four second-tier rows add a translation-level, a
transcript-level, an imaging and a knockdown task once the core four have been scored.

\pfstep{Essentiality, Goodall 2018} One TraDIS call per gene in each of
\goodallEnv{} media, \goodallGenes{} genes and \goodallLabels{} calls.
It is the simplest gene-level label in the host, and the one with three independent
calls to agree with, from the Keio construction outcome (Baba 2006), genetic
footprinting (Gerdes 2003) and the TnSeq no-insertion list (Price 2018). Genes are
keyed by \goodallIdentifier{}.

\pfstep{Chemical genomics, Nichols 2011} A colony-size S-score for every Keio deletion
in \nicholsEnv{} conditions: \nicholsGenes{} genes, \nicholsLabels{} scores. It is the
widest condition panel on one gene set in the host, so a held-out gene's whole profile
is the target. The RB-TnSeq compendium of Price 2018 poses the same task with a
sequencing readout on \priceEnv{} conditions and is the replication arm. Genes are keyed by
\nicholsIdentifier{}.

\pfstep{Proteome, Schmidt 2016} Absolute copies per cell for \schmidtGenes{} proteins
in \schmidtEnv{} growth conditions, \schmidtLabels{} values, on an unperturbed wild
type. The input is the coding sequence of the protein being measured, so this is the
row where a codon model is expected, untested here, to separate most from a one-hot
control. Mori 2021 measures the same quantity on a different platform and seven
samples; Caglar 2017 is the REL606 panel. Genes are keyed by \schmidtIdentifier{}.

\pfstep{Gene interaction, Babu 2014} A signed S score per deletion pair,
\babuLabels{} pairs over \babuGenes{} genes. It is the only pairwise label at scale in
the host and asks whether two gene vectors compose, which no single-gene task does.
Two splits are worth scoring: both pair members unseen, and one member unseen.
Butland 2008 is the unfiltered matrix of the same screens. Genes are keyed by
\babuIdentifier{}.

\pfstep{Second tier} Li 2014 gives an absolute protein synthesis rate per gene from
ribosome profiling in \liEnv{} media (\liGenes{} genes), the phenotype closest to what
a codon model reads. PRECISE-1K (Lamoureux 2023) gives a TPM per gene in \preciseEnv{}
distinct environments over \preciseInstances{} libraries (\preciseGenes{} genes), the
transcript-level counterpart of the proteome pick. Campos 2018 gives a vector of
cell-shape and cell-cycle features per Keio deletion (\camposGenes{} genes), the one
imaging phenotype in the host. Wang 2018 gives a gene-level CRISPRi fitness per gene
and screen (\wangGenes{} genes, \wangLabels{} values), the knockdown counterpart of
the deletion screens, and the one that reaches essential genes.

%% GENERATED FILE -- do not hand-edit.
%% SOURCE: experiments/database/scripts/build_ecoli_genslm_benchmark_table.py
%% off experiments/database/results/pre-build/2026-10-21/ecoli_datasets.json
\begin{landscape}
\begingroup
\scriptsize
\setlength{\tabcolsep}{3pt}
\begin{longtable}{@{}r >{\raggedright\arraybackslash}p{25mm} >{\raggedright\arraybackslash}p{20mm} >{\raggedright\arraybackslash}p{29mm} r r r >{\raggedright\arraybackslash}p{25mm} >{\raggedright\arraybackslash}p{23mm} >{\raggedright\arraybackslash}p{29mm} >{\raggedright\arraybackslash}p{46mm}@{}}
\caption{The short list for a gene-sequence encoding benchmark. \emph{Genes}, \emph{Labels}, \emph{Env} and \emph{Gene identifier} are read from a full scan of each store: distinct gene ids carrying a label, gene-attached numeric or boolean values, distinct stored environments, and the namespace the store keys genes by. \emph{Task}, \emph{Label}, \emph{Split}, \emph{Same task} and \emph{Why} are the pick list in the generating script.}\label{tab:shortlist}\\
\toprule
\# & Dataset & Task & Label & Genes & Labels & Env & Gene identifier & Split & Same task & Why \\
\midrule
\endfirsthead
\multicolumn{11}{@{}l}{\footnotesize\emph{Table~\ref{tab:shortlist}, continued}}\\
\toprule
\# & Dataset & Task & Label & Genes & Labels & Env & Gene identifier & Split & Same task & Why \\
\midrule
\endhead
\multicolumn{11}{@{}l}{\textbf{Core}} \\
1 & Gene\allowbreak{}Essentiality\allowbreak{}Goodall2018 & binary classification & essential or not, one TraDIS call per gene and medium & 4,181 & 8,309 & 2 & BW25113 locus tag, released as gene symbol & held-out genes & Gene\allowbreak{}Essentiality\allowbreak{}Baba2006, Gene\allowbreak{}Essentiality\allowbreak{}Gerdes2003, Gene\allowbreak{}Essentiality\allowbreak{}Price2018\allowbreak{}Ecoli & The simplest gene-level label in the host, and the one with three independent calls to agree with (Keio construction, genetic footprinting, TnSeq). A classifier on one vector per gene is the whole model. \\
2 & Env\allowbreak{}Chemgen\allowbreak{}Nichols2011 & regression, gene x condition & a colony-size S-score per (deletion, condition) & 3,949 & 1,286,962 & 324 & BW25113 locus tag, released as ECK accession & held-out genes, every condition & Rb\allowbreak{}Tnseq\allowbreak{}Price2018\allowbreak{}Ecoli, Env\allowbreak{}Chemgen\allowbreak{}Shiver2016, Carbon\allowbreak{}Source\allowbreak{}Tong2020, Env\allowbreak{}Chemgen\allowbreak{}Girgis2009 & The widest condition panel on one gene set. A held-out gene's whole profile is the target, and the condition enters as a second input the gene encoding does not have to carry. \\
3 & Proteome\allowbreak{}Schmidt2016 & regression, protein x condition & absolute copies per cell of one protein in one growth condition & 2,359 & 48,335 & 21 & gene symbol (readout key) & held-out genes, every condition & Proteome\allowbreak{}Mori2021, Proteome\allowbreak{}Caglar2017, Proteome\allowbreak{}Ishii2007 & A molecular level of the gene's own product, unperturbed, across conditions. Expected (untested) to be where a codon model separates from a one-hot control, because the input is the coding sequence whose product is measured. \\
4 & Gene\allowbreak{}Interaction\allowbreak{}Babu2014 & regression, gene pair & a signed colony-size S score per (query, array) deletion pair & 3,880 & 42,705 & 1 & MG1655 b-number, released as Keio JW number; b-number with a Keio sub-number; gene symbol & held-out genes (both pair members unseen), and one member unseen & Gene\allowbreak{}Interaction\allowbreak{}Butland2008, Crispri\allowbreak{}Cross\allowbreak{}Rachwalski2024 & The only pairwise label at scale in the host. It asks whether two gene vectors compose, which no single-gene task does. \\
\addlinespace
\multicolumn{11}{@{}l}{\textbf{Second tier}} \\
5 & Protein\allowbreak{}Synthesis\allowbreak{}Rate\allowbreak{}Li2014 & regression, gene x medium & absolute protein synthesis rate per gene from ribosome profiling & 3,410 & 8,644 & 3 & gene symbol (readout key) & held-out genes & Protein\allowbreak{}Turnover\allowbreak{}Gupta2024 & The phenotype closest to what a codon model reads: translation of the coding sequence itself, in three media. \\
6 & Rnaseq\allowbreak{}Lamoureux2023 & regression, gene x library & TPM of one gene in one RNA-seq library & 4,257 & 4,405,995 & 119 & b-number (readout key) & held-out genes, every library & Rnaseq\allowbreak{}Public\allowbreak{}K12\allowbreak{}Lamoureux2023, Rnaseq\allowbreak{}Caglar2017 & The transcript-level counterpart of the proteome pick, over a thousand libraries, so the per-gene target is a long expression profile. \\
7 & Morphology\allowbreak{}Campos2018 & multi-output regression & a vector of cell-shape and cell-cycle features per deletion & 3,664 & 53,508 & 1 & BW25113 locus tag & held-out genes & Growth\allowbreak{}Rate\allowbreak{}Campos2018 & The one imaging phenotype in the host: a different label family from growth and abundance, on the same Keio gene set. \\
8 & Crispri\allowbreak{}Gene\allowbreak{}Fitness\allowbreak{}Wang2018 & regression, gene x screen & gene-level CRISPRi fitness per (gene, screen) & 4,223 & 20,158 & 5 & MG1655 b-number & held-out genes & Crispri\allowbreak{}Knockdown\allowbreak{}Cui2018, Crispri\allowbreak{}Chemgen\allowbreak{}Choe2025, Mismatch\allowbreak{}Crispri\allowbreak{}Fitness\allowbreak{}Hawkins2020 & The knockdown counterpart of the deletion screens, already aggregated to the gene by its authors, so it covers essential genes the deletion collections cannot. \\
\bottomrule
\end{longtable}
\endgroup
\end{landscape}


%% notes-tex/database/ecoli-datasets/sections/caveats.tex
%% [[experiments.database.ecoli-datasets]]
%% https://github.com/Mjvolk3/torchcell/tree/main/notes-tex/database/ecoli-datasets/sections/caveats.tex
%% SOURCE: hand-authored prose; every number quoted here is a macro from
%% tables/numbers.tex, generated by
%% experiments/database/scripts/build_ecoli_genslm_benchmark_table.py

\section{What a score on these does not show\secstatus{todo}}\label{sec:caveats}

\pfstep{Gene identifiers are the released ones} Each store keys its genes as its
release does, and the \emph{Gene identifier} column of Table~\ref{tab:shortlist} is
measured from the records: b-numbers in one store, BW25113 locus tags or ECK accessions
in another, bare symbols in a third. Resolving those to the coding sequences of the
pinned assembly is a modeling-time step, by design, and its coverage must be measured
before any score is read, because a gene that fails to resolve is silently absent from
both the training and the test fold. The ECK synonym is the published one-to-one key
between the two K-12 strains and is the route for the Keio-based rows.

\pfstep{Held-out genes can leak through paralogs} A paralog of a held-out gene in the
training fold gives a sequence model a near copy of the test input. A split that
clusters genes by sequence similarity before holding out clusters removes that route; a
random split does not, and a random-split score is an upper bound.

\pfstep{The pretraining corpus is not disjoint from the test set} GenSLM was
pretrained on gene sequences from BV-BRC, a database that holds \org{E. coli} K-12
assemblies, so the genes being scored are expected, unverified for any one gene, to be
in its training data. The vector can then carry gene identity rather than a property
of the sequence, which is exactly what the one-hot control measures: an encoding that
scores no better than a one-hot vector of the same width has memorized, not
generalized.

\pfstep{Guide-level stores are excluded on purpose} The \guideStoreCount{} CRISPRi
stores that score one guide at a time are the largest \org{E. coli} rows, and the gene
each guide targets does have a coding sequence. The label, though, depends on the 20-nt
protospacer, its position in the gene and its knockdown strength, none of which the
gene's vector sees. Aggregating guides to their gene is what Wang 2018 released, and
that is the row on the list.

\pfstep{A tier is not a result} The sequence-only fit tier and the pick list say what
a sequence-only model could be given. No encoding has been scored on any of these rows
here; every expectation in Sec.~\ref{sec:benchmark} is labeled as one.


\end{document}
